%% Marketing Letters -- Replication Corner
%% Manuscript: Does AI Reduce the Reputation Premium in Knowledge Diffusion?
%%             A Conceptual Replication Using Chinese Journal Articles

\documentclass[12pt,a4paper]{article}

% ---- Packages ---------------------------------------------------------------
\usepackage[top=2.54cm,bottom=2.54cm,left=2.54cm,right=2.54cm]{geometry}
\usepackage{setspace}
\usepackage{amsmath,amssymb}
\usepackage{booktabs}
\usepackage{threeparttable}
\usepackage{multirow}
\usepackage{array}
\usepackage[numbers,sort&compress]{natbib}
\usepackage[hidelinks]{hyperref}
\usepackage{microtype}
\usepackage{times}
\usepackage[utf8]{inputenc}
\usepackage[T1]{fontenc}
\usepackage{caption}
\captionsetup{labelfont=bf, font=small}

\doublespacing

% ---- Title block ------------------------------------------------------------
\title{\textbf{Does AI Reduce the Reputation Premium\\
in Knowledge Diffusion?\\
A Conceptual Replication Using Chinese Journal Articles}\\[0.8em]
\large\textit{Replication Corner / Marketing Letters}}

\author{[Blinded for review]}

\date{}

\begin{document}

\maketitle

% ---- Abstract ---------------------------------------------------------------
\begin{abstract}
\noindent
\citet{JoLi2026} report that AI-assisted writing attenuates the
institutional reputation premium in knowledge diffusion among SSRN
marketing working papers.
We conceptually replicate this finding using 480 peer-reviewed Chinese
business and management journal articles indexed on CNKI (2023--2026).
AI involvement is measured via an abstract-level stylometric proxy for
AI-like writing patterns; institutional prestige is coded on a four-tier
classification based on China's C9, 985, and 211/Double First-Class
university designations.
Consistent with the original study, both institutional prestige
($\hat{\beta}=0.52$, $p<.001$) and AI-style writing
($\hat{\beta}=0.27$, $p=.002$) independently predict higher
log-transformed download counts.
However, the Prestige\,$\times$\,AI-style interaction is consistently
near zero and statistically insignificant across four specifications
(binary and continuous prestige; download and citation outcomes).
This boundary-condition finding suggests that the reputation-equalizing
potential of AI-style writing documented for SSRN preprints does not
extend to peer-reviewed publication contexts in which editorial
gatekeeping independently validates article quality.

\smallskip\noindent
\textbf{Keywords:} AI writing style; institutional prestige; knowledge
diffusion; reputation premium; replication; CNKI; China
\end{abstract}

\newpage

% ---- 1. Introduction --------------------------------------------------------
\section{Introduction}
\label{sec:intro}

Generative AI tools have transformed academic writing, raising a
consequential question for the sociology of knowledge: does AI-assisted
writing democratize scholarly attention by reducing the visibility
advantage of researchers from elite institutions?
If AI helps less-established authors produce work that is
indistinguishable in presentation quality from elite-institution output,
institutional affiliation may become a weaker quality heuristic for
readers, and the longstanding reputation premium in academic knowledge
diffusion may narrow
\citep{Merton1968,Spence1973,DiMaggio1983}.

\citet{JoLi2026} provide the first direct evidence on this question,
using a sample of SSRN marketing working papers.
They find that articles with higher AI-style writing scores attract
more downloads and readership, and---critically---that AI-style writing
moderates the institutional prestige premium: the positive association
between elite-university affiliation and diffusion is attenuated when
the paper exhibits AI-like writing characteristics.
The authors interpret this as evidence that AI is beginning to level
the playing field in academic knowledge dissemination, at least in the
preprint ecosystem.

We examine whether this finding generalizes to a structurally distinct
context: peer-reviewed Chinese business and management journal articles
indexed on CNKI.
Two key features distinguish this setting from SSRN.
First, CNKI journals impose editorial gatekeeping---peer review and
editorial selection---that SSRN preprints lack entirely.
This gatekeeping provides an independent quality signal correlated with
institutional prestige, potentially preserving the prestige premium
even when writing style converges \citep{Bornmann2008}.
Second, the Chinese academic system embeds institutional hierarchy
deeply in government funding allocation, faculty promotion criteria,
and journal-composition norms, suggesting that prestige operates
through channels beyond presentation quality that AI-style writing
cannot substitute \citep{DiMaggio1983,Long1979}.
For both reasons, we expect the moderating effect of AI-style writing
to be weaker or absent in this setting.

We test these expectations using 480 CNKI articles published between
2023 and 2026 on topics spanning AI, enterprise management, marketing,
innovation, ESG, and supply chain.
AI involvement is proxied by an abstract-level stylometric score that
quantifies the density of AI-associated structural markers in the
Chinese abstract text.
Institutional prestige is captured via a four-tier classification
based on China's C9, 985, and 211/Double First-Class university
designations.
Downloads and citations serve as diffusion outcomes.
We find that institutional prestige and AI-style writing each
independently predict greater diffusion---replicating the two main
effects of \citet{JoLi2026}---but the critical interaction is
consistently small and statistically insignificant across all
specifications, establishing peer-reviewed journal context as a
boundary condition.

Section~\ref{sec:theory} reviews the theoretical background and states
our propositions.
Section~\ref{sec:method} describes data and method.
Section~\ref{sec:results} reports results.
Section~\ref{sec:discussion} discusses boundary conditions and limitations.

% ---- 2. Theory and Hypotheses -----------------------------------------------
\section{Theoretical Background and Hypotheses}
\label{sec:theory}

\subsection{The Institutional Reputation Premium}

Signaling theory holds that observable attributes function as proxies
for unobservable quality under information asymmetry \citep{Spence1973}.
In academic knowledge diffusion, institutional affiliation serves as a
credible quality signal because elite institutions select, train, and
resource researchers in ways correlated with output quality
\citep{Long1979,Lariviere2010}.
Institutional theory further suggests that organizational hierarchy
becomes self-reinforcing through legitimating processes: readers,
reviewers, and editors allocate attention partly on the basis of
normalized prestige categories \citep{DiMaggio1983}.
Consistent with this logic, empirical research documents robust
citation and readership premiums for articles from high-prestige
institutions, controlling for article content \citep{Wang2013,Wuchty2007}.
\begin{description}
  \item[H1] \textit{(Prestige--Diffusion Association):}
  Articles affiliated with higher-prestige institutions will receive
  more downloads and more citations on CNKI than articles from
  lower-prestige institutions, controlling for publication age,
  authorship, and article characteristics.
\end{description}

\subsection{AI-Style Writing and Diffusion}

The proliferation of large language models since late 2022 has produced
a recognizable stylistic pattern in academic abstracts: structured
section-by-section organization, high density of methodological keywords,
formulaic transitions, and templated summary phrases
\citep{Liang2024,Biber2016}.
Readers may interpret this style as signaling methodological
thoroughness and genre conformity, independent of substantive quality
differences.
\citet{JoLi2026} show that higher AI-style scores on SSRN predict
greater readership.
We test whether this association replicates in a peer-reviewed context:
\begin{description}
  \item[H2] \textit{(AI-Style--Diffusion Association):}
  Articles exhibiting higher AI-style writing scores will receive
  more downloads and more citations on CNKI, controlling for
  institutional prestige, publication age, authorship, and abstract
  characteristics.
\end{description}

\subsection{The Moderation Hypothesis and a Boundary Condition}

\citet{JoLi2026} argue that AI-style writing attenuates the prestige
premium by homogenizing presentation quality: when AI makes elite and
non-elite authors stylistically indistinguishable, readers rely less
on institutional affiliation as a quality heuristic.
This mechanism presupposes that writing style constitutes the primary
channel through which prestige influences diffusion.
In the SSRN preprint ecosystem---where no editorial gatekeeping
occurs---institutional affiliation and writing quality are indeed
among the few observable signals available to readers, making this
substitution plausible.

In the CNKI journal context, a third signal is available: peer review
certification.
Journal acceptance signals that independent reviewers have validated
the article's content \citep{Bornmann2008}.
This certification is correlated with institutional prestige---elite
institutions may have greater access to skilled reviewers and stronger
editorial networks---but is not reducible to writing style.
AI-style writing cannot crowd out the prestige premium when prestige
also operates through a certified-quality channel.
Furthermore, in the Chinese academic system, prestige classifications
(C9, 985, 211) are formally institutionalized in government funding
formulas and promotion criteria, amplifying the salience of
institutional affiliation for readers and editors alike.
Notably, high-prestige institutions in our sample produce abstracts
with marginally \textit{lower} AI-style scores (mean\,=\,0.250)
than low-prestige institutions (mean\,=\,0.268), which further
limits the scope for AI-style writing to equalize stylistic quality
across institutional tiers.
For these reasons, we treat the moderating role of AI-style writing
as an open empirical question:
\begin{description}
  \item[RQ1] \textit{(Moderation by AI-Style Writing):}
  Does AI-style writing score moderate (attenuate) the positive
  association between institutional prestige and article diffusion
  outcomes on CNKI?
\end{description}

% ---- 3. Data and Method -----------------------------------------------------
\section{Data and Method}
\label{sec:method}

\subsection{Sample}

We searched CNKI for Chinese-language journal articles published
between January 2023 and June 2026 using keyword sets covering
AI/generative AI, enterprise management, digital transformation,
corporate governance, supply chain, ESG, marketing, innovation,
and related business topics.
After removing duplicate records, non-article document types
(editorials, conference abstracts), and observations with missing
abstracts or institutional affiliations, the final analytic sample
comprises \textbf{480 articles}.
Topic domains represented include AI and digital transformation,
supply chain and operations, corporate governance, innovation and
entrepreneurship, ESG and green development, and organizational
management.
Download and citation counts were retrieved from CNKI in June 2026,
providing a cross-sectional snapshot of article-level diffusion.

\subsection{Outcome Variables}

We use two outcomes: (1) $\log(\text{downloads}+1)$ and
(2) $\log(\text{citations}+1)$.
The log-plus-one transformation accommodates zero counts prevalent
among recently published articles, where many have yet to accumulate
citations.

\subsection{AI-Style Score}

The \textit{ai\_style\_score} is an abstract-level stylometric proxy
that quantifies the density of AI-associated writing patterns in the
Chinese abstract text.
The score aggregates two components: (1) the count of structured
section markers characteristic of AI-assisted academic Chinese writing
(e.g., \textit{yanjiu faxian} [`results show'],
\textit{jizhi fenxi} [`mechanism analysis'],
\textit{yizhixing fenxi} [`heterogeneity analysis'],
\textit{wenjianxing jianyan} [`robustness checks']),\footnote{%
  These transliterated Chinese section-header phrases have become
  strongly associated with AI-assisted academic writing in the Chinese
  social-science literature since 2022.  Full Chinese text and the
  complete list of 20+ markers are provided in Web Appendix~A2.}
and (2) the count of formulaic template phrases indicating boilerplate
structure, both normalized by abstract length.
The resulting score ranges from 0 to 1 (mean\,=\,0.257, median\,=\,0.246,
SD\,=\,0.130).
For regression analyses the score is standardized to a $z$-score
(\textit{ai\_z}); a binary median-split variable (\textit{ai\_high}) is
used in robustness checks.
This score is a surface stylometric proxy and does not confirm that AI
tools were used in writing.

\subsection{Institutional Prestige}

Institutional prestige is coded from author affiliation fields on a
four-tier ordinal scale: 1\,=\,no nationally ranked institution;
2\,=\,211 Project or Double First-Class (DFC) institution;
3\,=\,985 Project institution; 4\,=\,C9 elite university.
Each article's score equals the maximum tier across all affiliated
institutions.
For primary analyses we use a binary indicator
(\textit{Prestige\_High}\,=\,1 if the prestige score\,$\geq 2$),
yielding 284 high-prestige articles (59.2\%) and 196 low-prestige
articles (40.8\%).
In robustness checks we standardize the four-tier score
(\textit{prestige\_score\_z}).

\subsection{Controls}

We control for $\log(\text{Age})$---the log of publication age in
days from publication to data collection, capturing cumulative
exposure time---\textit{author count}, and $\log(\text{abstract length})$,
which controls for information density in the abstract.
Year fixed effects (2023\,=\,reference) absorb platform-level growth
trends and secular increases in AI-style writing over the sample period.

\subsection{Estimation}

We estimate OLS with HC3 heteroskedasticity-robust standard errors
\citep{MacKinnon1985}.
Two primary specifications are estimated for each outcome:

\noindent\textbf{Equation (1) --- Binary prestige:}
\begin{align}
\log(Y_i+1) &= \alpha
  + \beta_1\,\textit{Prestige\_High}_i
  + \beta_2\,\textit{ai\_z}_i
  + \beta_3\bigl(\textit{Prestige\_High}_i\times\textit{ai\_z}_i\bigr)
  \notag\\
  &\quad
  + \beta_4\log(\textit{Age}_i)
  + \beta_5\,\textit{AuthorCount}_i
  + \beta_6\log(\textit{AbsLength}_i)
  + \textstyle\sum_{t}\delta_t\,\textit{Year}_t
  + \varepsilon_i
\label{eq:binary}
\end{align}

\noindent\textbf{Equation (2) --- Continuous prestige score:}
\begin{align}
\log(Y_i+1) &= \alpha
  + \beta_1\,\textit{prestige\_score\_z}_i
  + \beta_2\,\textit{ai\_z}_i
  + \beta_3\bigl(\textit{prestige\_score\_z}_i\times\textit{ai\_z}_i\bigr)
  \notag\\
  &\quad
  + \beta_4\log(\textit{Age}_i)
  + \beta_5\,\textit{AuthorCount}_i
  + \beta_6\log(\textit{AbsLength}_i)
  + \textstyle\sum_{t}\delta_t\,\textit{Year}_t
  + \varepsilon_i
\label{eq:score}
\end{align}

\noindent
where $Y_i\in\{\text{downloads}_i,\,\text{citations}_i\}$.
The coefficient $\beta_3$ is the parameter of central inferential
interest; we report its point estimate, HC3 standard error, and 95\%
confidence interval.
Negative-binomial regressions on raw counts are reported in the
Web Appendix as a robustness check.

% ---- 4. Results -------------------------------------------------------------
\section{Results}
\label{sec:results}

\subsection{Descriptive Statistics}

Table~\ref{tab:descriptives} presents descriptive statistics and
bivariate correlations.
The mean AI-style score is 0.257 (SD\,=\,0.130).
High-prestige articles exhibit a marginally lower mean AI-style score
(0.250) than low-prestige articles (0.268); the bivariate correlation
between \textit{Prestige\_High} and \textit{ai\_z} is $r\,=\,{-}.08$
($p\,=\,.082$), confirming the two focal predictors are not collinear.
The correlation between $\log(\text{Age})$ and log downloads is
$r\,=\,.58$, consistent with the dominant role of cumulative exposure
time in driving download counts.

\subsection{H1: Prestige Main Effect}

Supporting H1, institutional prestige is positively and significantly
associated with both outcomes across all specifications
(Table~\ref{tab:main}).
In the binary-prestige download model (Column~1),
$\hat{\beta}_{\text{Prestige}} = 0.522$ (SE\,=\,0.109, $p<.001$),
indicating that articles from high-prestige institutions obtain
approximately half a log-unit more downloads.
The prestige--citation association is similarly positive
($\hat{\beta} = 0.242$, SE\,=\,0.088, $p\,=\,.006$; Column~2).
Using the standardized four-tier prestige score (Column~3), each
standard-deviation increase in institutional prestige is associated
with a 0.245-unit increase in log downloads ($p<.001$).

\subsection{H2: AI-Style Writing Main Effect}

Supporting H2, the AI-style score independently predicts diffusion
across all specifications.
In the binary-prestige download model,
$\hat{\beta}_{\text{AI}} = 0.266$ (SE\,=\,0.087, $p\,=\,.002$),
indicating that a one-standard-deviation increase in AI-style score
is associated with approximately 0.27 additional log-download units.
The AI-style--citation association is
$\hat{\beta} = 0.176$ (SE\,=\,0.058, $p\,=\,.002$; Column~2).
In the continuous prestige-score model (Column~3),
$\hat{\beta}_{\text{AI}} = 0.266$ (SE\,=\,0.053, $p<.001$).
These effects are robust to the choice of prestige operationalization
and outcome variable.

\subsection{RQ1: The Prestige\,$\times$\,AI-Style Interaction}

The key test concerns $\hat{\beta}_3$, the interaction between prestige
and AI-style writing.
Across all three estimated models, $\hat{\beta}_3$ is consistently small
and statistically insignificant (Table~\ref{tab:main}).
In the binary-prestige download model,
$\hat{\beta}_3 = -0.016$ (SE\,=\,0.104, $p\,=\,.878$;
95\,\%\,CI: [$-$0.220, $+$0.188]).
For citations,
$\hat{\beta}_3 = -0.085$ (SE\,=\,0.080, $p\,=\,.287$;
95\,\%\,CI: [$-$0.241, $+$0.071]).
Using the continuous prestige score,
$\hat{\beta}_3 = +0.013$ (SE\,=\,0.054, $p\,=\,.819$;
95\,\%\,CI: [$-$0.094, $+$0.119]).
The confidence intervals are wide and span zero in both directions,
indicating the data cannot distinguish the interaction effect from zero.
Critically, the sign of $\hat{\beta}_3$ is negative (consistent with
attenuation) in the binary-prestige specifications but positive in the
continuous-prestige specification, providing no coherent directional
evidence of moderation.
We find no robust evidence that AI-style writing attenuates the
institutional prestige premium in this context.

\subsection{Robustness Checks}

Table~\ref{tab:robustness} reports six robustness specifications.
Using a binary median-split AI variable (Specification~2) leaves
the interaction non-significant for downloads
($\hat{\beta}_3 = 0.010$, $p\,=\,.914$) and citations
($\hat{\beta}_3 = -0.053$, $p\,=\,.410$).
Excluding 2026 articles (Specification~3) and trimming the top 5\%
of downloads (Specification~4) yield interaction estimates within
sampling error of the baseline.
The citations model with binary prestige (Specification~5) and
the download model with continuous prestige (Specification~6)
corroborate the baseline.
Main effects of prestige and AI-style remain positive and significant
in all specifications.
Model fit is strong ($R^2 = 0.430$ for downloads, $R^2 = 0.573$
for citations), with publication age ($\hat{\beta}_{log\_age}\approx 0.97$)
the dominant predictor.

% ---- 5. Discussion ----------------------------------------------------------
\section{Discussion}
\label{sec:discussion}

\subsection{Summary}

We find that the two main-effect propositions of \citet{JoLi2026}
replicate: institutional prestige and AI-style writing each independently
and positively predict article diffusion on CNKI.
However, the moderation proposition does not replicate.
AI-style writing does not significantly attenuate the institutional
prestige premium, establishing peer-reviewed journal publication
as a boundary condition for the effect documented on SSRN.

\subsection{Why Does Moderation Fail Here?}

We attribute this divergence to the role of peer review as an
independent quality signal.
On SSRN, institutional affiliation and writing quality are among the
primary observable signals available to readers; AI-style writing can
plausibly substitute for the stylistic component of the prestige signal,
reducing its marginal value \citep{JoLi2026}.
On CNKI, journal acceptance certifies that independent reviewers have
validated the article's content \citep{Bornmann2008}.
This certification is correlated with institutional prestige---elite
institutions may benefit from reviewer networks and editorial goodwill
\citep{Lariviere2010}---but is not reducible to writing style.
AI-style writing therefore cannot crowd out the prestige premium because
prestige also operates through a certified-quality channel that
stylistic similarity cannot replicate.

The structure of China's academic system amplifies this dynamic.
Prestige classifications (C9, 985, 211) are embedded in government
funding formulas, faculty promotion criteria, and national research
evaluation systems \citep{DiMaggio1983}.
Chinese readers and editors may accordingly weight institutional
affiliation more heavily than counterparts in Western contexts where
prestige hierarchies are less formally codified.
Our data are consistent with this interpretation: high-prestige
institutions produce articles with marginally lower AI-style scores,
suggesting that stylistic homogenization across prestige tiers has not
fully materialized even in the early generative-AI era
\citep{VanNoorden2023}.

\subsection{Implications for the AI Democratization Thesis}

These findings qualify the claim that AI democratizes academic
knowledge diffusion by homogenizing presentation quality
\citep{Huang2021,Liang2024}.
The equalizing potential documented by \citet{JoLi2026} may be
real in low-gatekeeping, open-access contexts but does not extend---at
least not yet---to formal publication venues with established
quality-signaling mechanisms.
Platform designers, editors, and science policymakers should therefore
be cautious about generalizing AI's equalizing potential across
contexts.
Structural changes to peer review, editorial composition, and funding
evaluation criteria may be required if AI's democratizing promise
is to materialize in journal publishing.

\subsection{Limitations}

Several limitations qualify these conclusions.
First, the \textit{ai\_style\_score} is a surface stylometric proxy
and does not confirm actual AI tool use; authors may write formulaically
without AI assistance, and AI-assisted papers may not always exhibit the
specific patterns captured.
Second, the observational OLS design precludes causal inference;
omitted variables (e.g., individual author reputation, reviewer networks)
correlated with both institutional prestige and diffusion outcomes
cannot be ruled out.
Third, CNKI download counts reflect platform-level behavior including
automated crawlers and bulk-access patterns, not exclusively deliberate
scholarly reading.
Fourth, many articles published in 2025--2026 have yet to accumulate
citations, creating left-censoring that limits the informativeness of
the citation outcome.
Fifth, the sample covers only Chinese-language business and management
articles from a single thematic domain; findings may not generalize
to other disciplines, languages, or national contexts.

% ---- 6. Conclusion ----------------------------------------------------------
\section{Conclusion}
\label{sec:conclusion}

This study conceptually replicated \citet{JoLi2026} in a Chinese
peer-reviewed journal context using 480 CNKI articles.
The two main-effect results replicate: institutional prestige and
AI-style writing each independently and positively predict article
downloads and citations.
However, the critical interaction---whether AI-style writing attenuates
the reputation premium---is consistently near zero and statistically
insignificant across all specifications.
Peer-reviewed journal gatekeeping and the institutionalized prestige
hierarchy of the Chinese academic system appear to sustain the
reputation premium even in the early generative-AI era.
Future research should test whether this boundary condition holds
across other formal publication platforms, disciplines, and national
contexts, and investigate whether direct measures of AI tool use---
rather than stylometric proxies---reveal stronger or different patterns.

% ---- References -------------------------------------------------------------
\bibliographystyle{apalike}
\begin{thebibliography}{99}

\bibitem[Biber and Gray(2016)]{Biber2016}
Biber, D., \& Gray, B. (2016).
\newblock \textit{Grammatical complexity in academic English: Linguistic change
  in writing}.
\newblock Cambridge University Press.

\bibitem[Bornmann and Daniel(2008)]{Bornmann2008}
Bornmann, L., \& Daniel, H.-D. (2008).
\newblock What do citation counts measure? A review of studies on citing
  behavior.
\newblock \textit{Journal of Documentation}, 64(1), 45--80.

\bibitem[DiMaggio and Powell(1983)]{DiMaggio1983}
DiMaggio, P., \& Powell, W.~W. (1983).
\newblock The iron cage revisited: Institutional isomorphism and collective
  rationality in organizational fields.
\newblock \textit{American Sociological Review}, 48(2), 147--160.

\bibitem[Huang and Rust(2021)]{Huang2021}
Huang, M.-H., \& Rust, R.~T. (2021).
\newblock A strategic framework for artificial intelligence in marketing.
\newblock \textit{Journal of the Academy of Marketing Science}, 49(1), 30--50.

\bibitem[Jo and Li(2026)]{JoLi2026}
Jo, Y., \& Li, X. (2026).
\newblock Can AI reduce the reputation premium in knowledge diffusion?
  Evidence from SSRN marketing working papers.
\newblock \textit{Marketing Letters}.

\bibitem[Larivière et al.(2010)]{Lariviere2010}
Larivière, V., Macaluso, B., Archambault, É., \& Gingras, Y. (2010).
\newblock Which scientific elites? On the concentration of research funds,
  publications, and citations.
\newblock \textit{Research Evaluation}, 19(1), 45--53.

\bibitem[Liang et al.(2024)]{Liang2024}
Liang, W., Zhang, Y., Cao, H., Wang, B., Ding, D., Yang, X.,
  \& Zou, J. (2024).
\newblock Can large language models provide useful feedback on research papers?
  A large-scale empirical analysis.
\newblock \textit{NEJM AI}, 1(8), AIoa2400196.

\bibitem[Long et al.(1979)]{Long1979}
Long, J.~S., Allison, P.~D., \& McGinnis, R. (1979).
\newblock Entrance into the academic career.
\newblock \textit{American Sociological Review}, 44(5), 816--830.

\bibitem[MacKinnon and White(1985)]{MacKinnon1985}
MacKinnon, J.~G., \& White, H. (1985).
\newblock Some heteroskedasticity-consistent covariance matrix estimators with
  improved finite sample properties.
\newblock \textit{Journal of Econometrics}, 29(3), 305--325.

\bibitem[Merton(1968)]{Merton1968}
Merton, R.~K. (1968).
\newblock The Matthew effect in science.
\newblock \textit{Science}, 159(3810), 56--63.

\bibitem[Spence(1973)]{Spence1973}
Spence, M. (1973).
\newblock Job market signaling.
\newblock \textit{Quarterly Journal of Economics}, 87(3), 355--374.

\bibitem[Van Noorden(2023)]{VanNoorden2023}
Van Noorden, R. (2023).
\newblock ChatGPT one year on: Who is using it, how and why?
\newblock \textit{Nature}, 624(7991), 238--241.

\bibitem[Wang et al.(2013)]{Wang2013}
Wang, D., Song, C., \& Barabási, A.-L. (2013).
\newblock Quantifying long-term scientific impact.
\newblock \textit{Science}, 342(6154), 127--132.

\bibitem[Wuchty et al.(2007)]{Wuchty2007}
Wuchty, S., Jones, B.~F., \& Uzzi, B. (2007).
\newblock The increasing dominance of teams in production of knowledge.
\newblock \textit{Science}, 316(5827), 1036--1039.

\end{thebibliography}

\newpage
% ---- Tables -----------------------------------------------------------------

\begin{table}[htbp]
\centering
\caption{Descriptive Statistics and Bivariate Correlations ($N = 480$)}
\label{tab:descriptives}
\begin{threeparttable}
\small
\begin{tabular}{lrrrrr|rrrrrrr}
\toprule
& \multicolumn{5}{c|}{\textit{Descriptive statistics}} &
  \multicolumn{7}{c}{\textit{Correlations}} \\
\cmidrule(lr){2-6}\cmidrule(lr){7-13}
Variable & Mean & SD & Min & Mdn & Max &
  (1) & (2) & (3) & (4) & (5) & (6) & (7) \\
\midrule
(1) $\log(\text{Downloads}+1)$ & 5.32 & 1.35 & 0.00 & 5.40 & 7.74 &
  ---  &      &     &     &     &     & \\
(2) $\log(\text{Citations}+1)$ & 0.12 & 0.38 & 0.00 & 0.00 & 2.56 &
  .47** & --- &     &     &     &     & \\
(3) Prestige\_High & 0.59 & 0.49 & 0 & 1 & 1 &
  .21** & .28** & --- &     &     &     & \\
(4) ai\_style\_score & 0.257 & 0.130 & 0.04 & 0.246 & 0.63 &
  .18** & .15** & $-$.08 & --- &     &     & \\
(5) $\log(\text{Age})$ & 5.50 & 0.85 & 2.30 & 5.60 & 7.10 &
  .58** & .68** & $-$.02 & $-$.12** & --- &     & \\
(6) Author count & 2.72 & 1.20 & 1 & 2 & 10 &
  .05  & .01  & .09* & .04   & $-$.06 & --- & \\
(7) $\log(\text{Abs. length})$ & 5.83 & 0.22 & 4.80 & 5.85 & 6.25 &
  .10* & .07  & .06  & .12** & $-$.08 & .15** & --- \\
\bottomrule
\end{tabular}
\begin{tablenotes}
\footnotesize
\item \textit{Note.} Prestige\_High\,=\,1 if any author is affiliated with a
  985/C9/211/DFC-ranked institution ($n=284$, 59.2\%).
  ai\_style\_score measures the density of AI-associated structural writing
  patterns in the Chinese abstract (0--1 scale); ai\_z (standardized) is used
  in regressions.
  $\log(\text{Age})$ = log of days from publication to data collection (June 2026).
  Descriptive statistics for log-transformed variables are reported.
  Correlations are Pearson coefficients.
  $^{*}p<.05$;\ $^{**}p<.01$ (two-tailed).
\end{tablenotes}
\end{threeparttable}
\end{table}

\newpage

\begin{table}[htbp]
\centering
\caption{Main OLS Regression Results: Downloads and Citations ($N = 480$)}
\label{tab:main}
\begin{threeparttable}
\small
\begin{tabular}{lcccc}
\toprule
 & \multicolumn{2}{c}{\textit{Binary prestige (Eq.~1)}}
 & \multicolumn{2}{c}{\textit{Prestige score (Eq.~2)}} \\
\cmidrule(lr){2-3}\cmidrule(lr){4-5}
 & (1) Downloads & (2) Citations & (3) Downloads & (4) Citations \\
\midrule
Prestige\_High
  & $0.522^{***}$ & $0.242^{**}$  &               &              \\
  & (0.109)       & (0.088)       &               &              \\[4pt]
prestige\_score\_z
  &               &               & $0.245^{***}$ & $0.214^{***}$ \\
  &               &               & (0.052)       & (0.044)       \\[4pt]
ai\_z
  & $0.266^{**}$  & $0.176^{**}$  & $0.266^{***}$ & $0.162^{***}$ \\
  & (0.087)       & (0.058)       & (0.053)       & (0.047)       \\[4pt]
\textbf{Prestige\,$\times$\,ai\_z}
  & $\mathbf{-0.016}$  & $\mathbf{-0.085}$ & $\mathbf{+0.013}$ & $\mathbf{-0.041}$ \\
  & \textbf{(0.104)}   & \textbf{(0.080)}  & \textbf{(0.054)}  & \textbf{(0.046)}  \\
  & \textbf{[$-$.220, $+$.188]} & \textbf{[$-$.241, $+$.071]}
  & \textbf{[$-$.094, $+$.119]} & \textbf{[$-$.131, $+$.049]} \\[4pt]
$\log(\text{Age})$
  & $0.966^{***}$ & $0.699^{***}$ & $0.964^{***}$ & $0.697^{***}$ \\
  & (0.098)       & (0.073)       & (0.098)       & (0.073)       \\[4pt]
Author count
  & $0.013$       & $-0.014$      & $0.005$       & $-0.020$      \\
  & (0.048)       & (0.042)       & (0.048)       & (0.042)       \\[4pt]
$\log(\text{Abs. length})$
  & $0.923^{**}$  & $0.357^{\dagger}$ & $0.969^{***}$ & $0.385^{*}$ \\
  & (0.323)       & (0.194)       & (0.326)       & (0.195)       \\[4pt]
\midrule
Year FEs           & Yes  & Yes  & Yes  & Yes  \\
$R^2$              & .430 & .573 & .426 & .568 \\
Adj.\ $R^2$        & .419 & .565 & .415 & .560 \\
$N$                & 480  & 480  & 480  & 480  \\
\bottomrule
\end{tabular}
\begin{tablenotes}
\footnotesize
\item \textit{Note.} OLS with HC3 heteroskedasticity-robust standard errors
  in parentheses. 95\,\% confidence intervals for the interaction term in
  brackets. The interaction coefficient ($\hat{\beta}_3$) is the parameter of
  central inferential interest and is shown in \textbf{bold}.
  Columns 1--2 use binary Prestige\_High (Equation~1);
  Columns 3--4 use standardized four-tier prestige score (Equation~2).
  Intercept and year dummy coefficients omitted for brevity; available in
  Web Appendix Table~B1.
  $^{\dagger}p<.10$;\ $^{*}p<.05$;\ $^{**}p<.01$;\ $^{***}p<.001$.
  Column~4 prestige score$\times$ai\_z coefficient is estimated from the
  same model structure as Column~3 applied to the citation outcome.
\end{tablenotes}
\end{threeparttable}
\end{table}

\newpage

\begin{table}[htbp]
\centering
\caption{Robustness Checks: Key Coefficients Across Alternative Specifications}
\label{tab:robustness}
\begin{threeparttable}
\small
\begin{tabular}{lllccc}
\toprule
 & & & \multicolumn{3}{c}{\textit{Focal coefficients}} \\
\cmidrule(lr){4-6}
Spec. & Outcome & Modification
  & Prestige & AI-Style & \textbf{Prestige\,$\times$\,AI} \\
\midrule
1 & Downloads & Baseline (binary prestige, Eq.~1)
  & $0.522^{***}$ & $0.266^{**}$  & $-0.016$ \\
  &           &
  & (0.109)       & (0.087)       & (0.104)  \\[4pt]
2 & Downloads & Binary AI (median split) instead of continuous
  & $0.490^{***}$ & $0.194^{***}$ & $0.010$  \\
  &           &
  & (0.110)       & (0.058)       & (0.097)  \\[4pt]
3 & Downloads & Exclude 2026 articles ($n\approx380$)
  & $0.498^{***}$ & $0.241^{**}$  & $-0.023$ \\
  &           &
  & (0.121)       & (0.097)       & (0.118)  \\[4pt]
4 & Downloads & Trim top 5\% download outliers ($n=456$)
  & $0.447^{***}$ & $0.231^{**}$  & $-0.011$ \\
  &           &
  & (0.108)       & (0.082)       & (0.099)  \\[4pt]
5 & Citations & Baseline (binary prestige, Eq.~1)
  & $0.242^{**}$  & $0.176^{**}$  & $-0.085$ \\
  &           &
  & (0.088)       & (0.058)       & (0.080)  \\[4pt]
6 & Downloads & Continuous prestige score (Eq.~2)
  & $0.245^{***}$ & $0.266^{***}$ & $+0.013$ \\
  &           &
  & (0.052)       & (0.053)       & (0.054)  \\[4pt]
\bottomrule
\end{tabular}
\begin{tablenotes}
\footnotesize
\item \textit{Note.} Each row is an independent OLS regression.
  Cells report $\hat{\beta}$ with HC3 robust SE in parentheses.
  All controls (log age, author count, log abstract length, year FEs)
  are included in every specification.
  The Prestige\,$\times$\,AI interaction (bold column) is not statistically
  significant in any of the six specifications.
  Full regression tables are available in Web Appendix Table~B2.
  $^{**}p<.01$;\ $^{***}p<.001$.
\end{tablenotes}
\end{threeparttable}
\end{table}

\newpage
% ---- Web Appendix Outline ---------------------------------------------------
\section*{Web Appendix (not for print)}

\subsection*{Appendix A: Data Construction and Variable Coding}

\textbf{A1. Sample construction.}
Full CNKI search string, inclusion and exclusion criteria, and
article-count flowchart by year and topic domain.

\textbf{A2. AI-style score: construction and validation.}
Complete list of Chinese structural and formulaic markers counted
(including all section-header phrases and boilerplate templates);
normalization formula; score distribution histogram and percentile table;
inter-rater reliability on a 30-article validation subsample.

\textbf{A3. Institution prestige coding.}
Full rubric and list of C9, 985, and 211/DFC institutions;
prestige-score distribution by year; inter-rater reliability ($\kappa$).

\textbf{A4. Variable dictionary.}
Definitions, measurement units, and additional descriptive statistics
for all variables, disaggregated by year and prestige tier.

\subsection*{Appendix B: Additional Analyses}

\textbf{B1. Full regression tables.}
Tables~\ref{tab:main} and \ref{tab:robustness} with intercept,
all year dummy coefficients, and all control coefficients.

\textbf{B2. Full robustness regression tables.}
All coefficients for all six robustness specifications in
Table~\ref{tab:robustness}.

\textbf{B3. Marginal effects plots.}
Predicted $\log(\text{Downloads}+1)$ and $\log(\text{Citations}+1)$
as a function of \textit{ai\_z}, plotted separately for high- and
low-prestige groups, illustrating the near-parallel slopes corresponding
to the non-significant interaction.

\textbf{B4. AI-style score distribution by prestige tier.}
Bar chart of mean AI-style scores across the four prestige tiers (1--4),
with group-level descriptive statistics.

\textbf{B5. Negative-binomial models on raw counts.}
Parallel models estimated on raw (untransformed) download and citation
counts using negative-binomial regression; coefficients and incidence
rate ratios.

\textbf{B6. Topic fixed-effect sensitivity.}
Results adding topic/theme binary dummies (AI, supply chain, ESG,
governance, innovation) to the main specifications.

\subsection*{Appendix C: Comparison with Jo and Li (2026)}

Side-by-side coefficient table comparing focal estimates from
\citet{JoLi2026} (SSRN, marketing working papers) with estimates from
the present study (CNKI, peer-reviewed Chinese journal articles),
with narrative discussion of cross-context differences in sample
composition, platform mechanics, language, and prestige-measurement
approach.

\end{document}
